\chapter{Kiểm tra và làm sạch dữ liệu}\label{ch:lamsach}

\section{Kiểm tra ban đầu}\label{sec:kiemtra}
Dữ liệu thô được kiểm tra bằng hàm \texttt{audit\_data()} của \texttt{pipeline.py} (Bảng~\ref{tab:kiemtra}, \ref{tab:loi}).

\begin{table}[H]
\centering
\caption{Kết quả kiểm tra ban đầu trên tệp gốc (13,580 dòng)}\label{tab:kiemtra}
\small
\begin{tabular}{@{}l L{10.3cm}@{}}
\toprule
Khía cạnh & Kết quả \\
\midrule
Kiểu dữ liệu & 8 cột chuỗi, 1 số nguyên, 12 số thực; \texttt{Date} là chuỗi \texttt{dd/mm/yyyy}, đọc bằng \texttt{dayfirst=True} \\
Dữ liệu thiếu & \texttt{BuildingArea} 6,450 dòng (47.5\%), \texttt{YearBuilt} 5,375 (39.6\%), \texttt{CouncilArea} 1,369 (10.1\%), \texttt{Car} 62 (0.5\%); 17 cột còn lại đầy đủ \\
Trùng lặp & 0 dòng trùng hoàn toàn; 13,378 địa chỉ khác nhau trên 13,580 dòng nên một số địa chỉ lặp lại, nhưng không dòng nào trùng toàn bộ nên đều được giữ \\
Giá trị lỗi hoặc bất thường & Bảng~\ref{tab:loi} \\
\bottomrule
\end{tabular}
\nguon{\texttt{outputs/data\_quality\_profile.csv} (kiểu, số ô thiếu, \% thiếu); số dòng trùng và địa chỉ: \bs{N03, N07, N08}.}
\end{table}

\begin{table}[H]
\centering
\caption{Giá trị lỗi hoặc bất thường trên tệp gốc và cách xử lý}\label{tab:loi}
\small
\begin{tabular}{@{}L{5.4cm} r L{5.6cm}@{}}
\toprule
Điều kiện & Số dòng & Xử lý \\
\midrule
\texttt{Price} $\leq$ 0, \texttt{Rooms} $\leq$ 0, \texttt{Distance} $<$ 0, \texttt{Landsize} $<$ 0 & 0 & Có kiểm tra, không dòng nào vi phạm \\
Tọa độ (\texttt{Lattitude}, \texttt{Longtitude}) ngoài khung Melbourne & 0 & Có kiểm tra, không dòng nào vi phạm \\
\texttt{BuildingArea} $\leq$ 0 & 17 & Chuyển thành thiếu \\
\texttt{YearBuilt} $<$ 1700 & 1 & Chuyển thành thiếu; giá trị nhỏ nhất là năm 1196 \\
\texttt{Landsize} $>$ 10,000 & 26 & Chuyển thành thiếu; lớn nhất 433,014 m\textsuperscript{2}, 15 dòng là căn hộ \\
\texttt{BuildingArea} $>$ 1,000 & 8 & Chuyển thành thiếu; lớn nhất 44,515 m\textsuperscript{2} \\
\texttt{Bathroom} $=$ 0 & 34 & Chuyển thành thiếu \\
\texttt{Landsize} $=$ 0 & 1,939 & Giữ nguyên; 1,629 dòng (84\%) là căn hộ \\
\texttt{YearBuilt} $\geq$ 2017 & 19 & Giữ nguyên; 6 dòng có năm xây sau năm bán \\
\texttt{Distance} $=$ 0; \texttt{Rooms} $\geq$ 7 & 6; 19 & Giữ nguyên (nhà ngay trung tâm, nhà lớn) \\
\bottomrule
\end{tabular}
\nguon{\bs{L01, L02}; số dòng vi phạm do \texttt{audit\_data()} đếm, xem thêm \texttt{pipeline.py}.}
\end{table}

Khung tọa độ kiểm tra là \texttt{Lattitude} trong $[-39.5;\,-36.0]$ và \texttt{Longtitude} trong $[142.0;\,147.0]$. Trên tệp gốc \texttt{Lattitude} nằm trong $[-38.18;\,-37.41]$ và \texttt{Longtitude} trong $[144.43;\,145.53]$ (\bs{L01}).

\begin{figure}[H]
\centering
\includegraphics[width=0.78\linewidth]{f_missing.png}
\caption{Tỷ lệ dữ liệu thiếu theo cột trong tệp gốc}\label{fig:missing}
\nguon{\texttt{report/make\_figures.py}, hàm \texttt{f\_missing()}, từ \texttt{data/melb\_data.csv}.}
\end{figure}

\texttt{BuildingArea} và \texttt{YearBuilt} thiếu gần một nửa: loại mọi dòng thiếu một trong hai cột sau làm sạch sẽ mất 6,744 dòng (49.7\%) (\bs{L03}). Vì vậy dự án không xóa dòng mà xử lý giá trị thiếu trong \texttt{Pipeline} của mô hình (mục~\ref{sec:thieu}).

\section{Quy tắc làm sạch}\label{sec:quytac}
Hàm \texttt{clean\_data()} của \texttt{pipeline.py} chuyển giá trị không hợp lệ thành \texttt{NaN} thay vì xóa dòng và không dùng thống kê của tập dữ liệu. Giá trị \emph{chắc chắn không hợp lệ hoặc phi thực tế cho một căn nhà} (diện tích xây bằng 0 hoặc trên 1,000 m\textsuperscript{2}, đất trên 10,000 m\textsuperscript{2}, phòng tắm bằng 0, năm xây trước 1700) bị chuyển thành thiếu. Giá trị đáng ngờ nhưng có thể có thật (đất bằng 0 ở căn hộ, giá bán cao) được giữ nguyên. Hai ngưỡng 10,000 và 1,000 m\textsuperscript{2} là lựa chọn của dự án (hằng số \texttt{MAX\_LANDSIZE}, \texttt{MAX\_BUILDING\_AREA}), đặt sau khi xem các dòng vượt ngưỡng; tác động của chúng được đánh giá ở Phép 6 và mục~\ref{sec:luuy}. Bảng~\ref{tab:quytac} liệt kê quy tắc; mã nằm ở Phụ lục~\ref{pl:quytac}.

\begin{table}[H]
\centering
\caption{Quy tắc của \texttt{clean\_data()}, lý do và số dòng bị ảnh hưởng}\label{tab:quytac}
\small
\begin{tabular}{@{}c L{4.2cm} L{5.6cm} r@{}}
\toprule
\# & Quy tắc & Lý do & Số dòng \\
\midrule
1 & Loại dòng trùng hoàn toàn & Tránh đếm một giao dịch hai lần & 0 \\
2 & \texttt{Date} $\to$ kiểu ngày, \texttt{dayfirst=True} & Dữ liệu ghi ngày/tháng/năm & 13,580 \\
3 & Ép các cột số về kiểu số & Giá trị không chuyển được sẽ thành thiếu & 0 \\
4 & Cờ \texttt{*\_was\_missing} cho \texttt{Car}, \texttt{BuildingArea}, \texttt{YearBuilt}, \texttt{CouncilArea} & Ghi nhận ô vốn đã thiếu \emph{trước} khi các quy tắc 5--8 đặt thêm giá trị thiếu & 62; 6,450; 5,375; 1,369 \\
5 & \texttt{BuildingArea} $\leq$ 0 hoặc $>$ 1,000 $\to$ thiếu & Bằng 0 hoặc cỡ khu phức hợp là lỗi nhập liệu & 17; 8 \\
6 & \texttt{YearBuilt} $<$ 1700 $\to$ thiếu & Năm 1196 là lỗi nhập liệu & 1 \\
7 & \texttt{Landsize} $>$ 10,000 $\to$ thiếu & Đất trên 1 ha không đại diện cho một căn nhà & 26 \\
8 & \texttt{Bathroom} $=$ 0 $\to$ thiếu & Nhà ở phải có phòng tắm & 34 \\
9 & \texttt{Price}, \texttt{Rooms} $\leq$ 0; \texttt{Distance}, \texttt{Landsize} $<$ 0; tọa độ ngoài khung $\to$ thiếu & Giá trị bất khả thi & 0 \\
10 & Loại dòng không có giá hợp lệ & Giá là biến mục tiêu, không điền & 0 \\
11 & \texttt{CouncilArea} thiếu $\to$ ``Unknown'' & Giữ ``không biết'' thành nhóm riêng & 1,369 \\
12 & Thêm \texttt{SaleYear}, \texttt{SaleMonth} & Phục vụ phân tích theo thời gian và làm đặc trưng & 13,580 \\
\bottomrule
\end{tabular}
\nguon{\texttt{pipeline.py}, hàm \texttt{clean\_data()}; số dòng: \bs{L01, L02}.}
\end{table}

\begin{figure}[H]
\centering
\includegraphics[width=0.78\linewidth]{f_cleaning.png}
\caption{Số dòng bị ảnh hưởng bởi từng quy tắc làm sạch}\label{fig:cleaning}
\nguon{\texttt{report/make\_figures.py}, hàm \texttt{f\_cleaning()}, đếm từ \texttt{data/melb\_data.csv}.}
\end{figure}

Không dòng nào bị xóa: bảng sạch giữ 13,580 dòng và 27 cột. Số ô thiếu sau làm sạch: \texttt{BuildingArea} 6,475 (6,450 cộng 25 ô quy tắc 5), \texttt{YearBuilt} 5,376, \texttt{Landsize} 26, \texttt{Bathroom} 34, \texttt{Car} 62, \texttt{CouncilArea} 0 (\bs{L03}).

\section{Dữ liệu thiếu và ngoại lai}\label{sec:thieu}
\paragraph{Dữ liệu thiếu.} \texttt{pipeline.py} không điền \texttt{BuildingArea} và \texttt{YearBuilt}, vì tỷ lệ thiếu cao và phép điền có thể làm méo thống kê mô tả; các phân tích dùng những dòng có giá trị. Việc điền chỉ thực hiện trong \texttt{Pipeline} của mô hình: cột số điền bằng trung vị của tập train, cột phân loại bằng giá trị phổ biến nhất của tập train. Các cờ \texttt{*\_was\_missing} được đưa vào mô hình vì việc thiếu không ngẫu nhiên: giá trung vị của dòng thiếu \texttt{BuildingArea} là 917,000 AUD so với 895,000 ở dòng có giá trị; với \texttt{YearBuilt} là 937,500 so với 880,000 (\bs{L05}). Các dòng \texttt{Landsize} bằng 0 có giá trung vị 560,000 so với 975,000 ở dòng có đất, vì 84\% trong số đó là căn hộ (\bs{L02, L05}).

\paragraph{Ngoại lai về giá.} Theo quy tắc tứ phân vị (Q1 $=$ 650,000; Q3 $=$ 1,330,000), ngưỡng trên là 2,350,000 AUD và có 612 dòng (4.5\%) vượt ngưỡng. Các dòng này được giữ vì là giao dịch có thể có thật; lớn nhất là một căn ba phòng ở Mulgrave giá 9,000,000 AUD, cách trung tâm 18.8 km (\bs{L04}). Để giảm ảnh hưởng của đuôi phải, Random Forest được huấn luyện trên $\log(1+\text{giá})$ (\texttt{train\_model.py}); độ lệch của giá giảm từ 2.24 xuống 0.18 khi lấy logarit (\bs{L04}, Hình~\ref{fig:hist}).

\paragraph{Ngoại lai về diện tích.} Khác với giá, giá trị diện tích cực đoan là lỗi nhập liệu hoặc lô đất của cả khu phức hợp, nên bị chuyển thành thiếu (quy tắc 5, 7) rồi được điền bằng trung vị tập train trong \texttt{Pipeline}. Việc này đưa Pearson giữa \texttt{BuildingArea} và giá từ 0.09 lên 0.56 (Phép 6). Dòng \texttt{Landsize} bằng 0 không bị coi là ngoại lai vì phần lớn là căn hộ.

\section{Chia tập train, validation, test}\label{sec:chia}
\texttt{train\_model.py} chia \emph{ngẫu nhiên} hai lần với \texttt{random\_state=42}: tách 20\% làm tập test, rồi tách 20\% phần còn lại làm tập validation. Tập test chỉ dùng để chấm điểm mô hình đã chọn.

\begin{table}[H]
\centering
\caption{Ba tập dữ liệu của \texttt{train\_model.py}}\label{tab:chia}
\small
\begin{tabular}{@{}l r r r@{}}
\toprule
Tập & Số dòng & Tỷ lệ & Giá trung vị (AUD) \\
\midrule
Train & 8,691 & 64\% & 908,000 \\
Validation & 2,173 & 16\% & 895,000 \\
Test & 2,716 & 20\% & 900,000 \\
\bottomrule
\end{tabular}
\nguon{\texttt{train\_model.py} (\texttt{train\_test\_split}); số dòng và trung vị: \bs{L06}.}
\end{table}

Giá trung vị của ba tập gần nhau (895,000 đến 908,000). Tuy nhiên dữ liệu có ngày bán và phép chia ngẫu nhiên để giao dịch \emph{sau} nằm lẫn trong tập train, nên sai số ở Chương~\ref{ch:mohinh} có thể lạc quan hơn so với dự đoán cho giai đoạn chưa xảy ra. Chia theo thời gian là hướng cải thiện (mục~\ref{sec:hanche}).

\paragraph{Chỉ học thống kê từ tập train.} Điền thiếu, chuẩn hóa (\texttt{StandardScaler}) và one-hot nằm trong \texttt{Pipeline} nên chỉ được \texttt{fit} trên tập train. Đoạn mã dưới rút gọn từ \texttt{build\_preprocessor()}.

\begin{lstlisting}
numeric_pipeline = Pipeline([("imputer", SimpleImputer(strategy="median")),
                             ("scaler", StandardScaler())])
categorical_pipeline = Pipeline([("imputer", SimpleImputer(strategy="most_frequent")),
                                 ("onehot", OneHotEncoder(handle_unknown="ignore"))])
ColumnTransformer([("numeric", numeric_pipeline, numeric_features),
                   ("categorical", categorical_pipeline, categorical_features)])
\end{lstlisting}

\texttt{handle\_unknown="ignore"} cho phép dự đoán khi gặp suburb chưa có trong tập train.

\section{Luồng xử lý giữa các chương trình}\label{sec:dungchung}
Dự án gồm hai bước nối tiếp, ghi kết quả trung gian ra tệp (Hình~\ref{fig:dungchung}).

\begin{figure}[H]
\centering
\begin{tikzpicture}[node distance=6mm, every node/.style={font=\footnotesize},
  box/.style={draw,rounded corners,align=center,minimum height=9mm,fill=blue!6},
  arr/.style={-{Stealth[length=2.2mm]}}]
  \node[box] (raw) {Tệp gốc\\\texttt{data/melb\_data.csv}};
  \node[box,right=of raw] (pipe) {\texttt{pipeline.py}\\làm sạch, phân tích};
  \node[box,right=of pipe] (clean) {\texttt{outputs/}\\\texttt{melb\_data\_clean.csv}};
  \node[box,below=10mm of clean,fill=green!8] (train) {\texttt{train\_model.py}\\chia tập, huấn luyện, đánh giá};
  \node[box,left=of train,fill=orange!10] (mod) {\texttt{models/}\\\texttt{best\_pipeline.joblib}};
  \draw[arr] (raw)--(pipe); \draw[arr] (pipe)--(clean); \draw[arr] (clean)--(train); \draw[arr] (train)--(mod);
\end{tikzpicture}
\caption{Luồng dữ liệu của dự án}\label{fig:dungchung}
\nguon{\texttt{README.md} và mã nguồn (\texttt{pipeline.py}, \texttt{train\_model.py}).}
\end{figure}

Notebook \texttt{melbourne\_data\_analysis\_pipeline.ipynb} lặp lại các quy tắc làm sạch và phân tích theo từng bước. Tệp \texttt{.joblib} phụ thuộc phiên bản thư viện, nên cần chạy lại \texttt{train\_model.py} trên máy sử dụng.
